\documentclass[../../main.tex]{subfiles}% 注意这里的文件路径不能用 ./main.tex ,否则用latexmk编译子文件会报错
\graphicspath{{\subfix{./image/}}} % 指定图片目录,后续可以直接使用图片文件名
% 注意这里的文件路径不能用 ../../image/ ,否则用latexmk编译子文件会报错

% 例如:
% \begin{figure}[H]
% \centering
% \includegraphics[scale=0.3]{图.png}
% \caption{}
% \label{figure:图}
% \end{figure}
% 注意:上述\label{}一定要放在\caption{}之后,否则引用图片序号会只会显示??.

\begin{document}

\section{重积分计算}

\begin{example}
设$f(x,y)$在$D$上有定义，且在某一圆盘外为$0$。求证：
\begin{align*}
\iint_D \left[ \frac{\partial^2 f}{\partial x^2} \frac{\partial^2 f}{\partial y^2} - \left( \frac{\partial^2 f}{\partial x \partial y} \right)^2 \right] \mathrm{d}x \mathrm{d}y = 0.
\end{align*}
\end{example}
\begin{proof}

由题设可知$f$在$\partial D$上恒为$0$，进而
$$\frac{\partial f}{\partial x}(x,y) = \frac{\partial f}{\partial y}(x,y) = 0, \quad \forall (x,y) \in \partial D.$$
由Green公式可得
\begin{align*}
&\iint_D \left[ \frac{\partial^2 f}{\partial x^2} \frac{\partial^2 f}{\partial y^2} - \left( \frac{\partial^2 f}{\partial x \partial y} \right)^2 \right] \mathrm{d}x \mathrm{d}y
= \iint_D (f_{xx} f_{yy} - f_{xy}^2) \mathrm{d}x \mathrm{d}y \\
&= \iint_D \left[ (f_{xx} f_y)_y - (f_{xy} f_y)_x \right] \mathrm{d}x \mathrm{d}y 
\xlongequal{\text{Green公式}} \oint_{\partial D} f_{xy} f_y \mathrm{d}x + f_{xx} f_y \mathrm{d}y = 0.
\end{align*}

\end{proof}

\begin{example}
设不全为0的实数$\alpha,\beta,\gamma$，考虑
\begin{align*}
C:
\begin{cases}
x^2 + y^2 + z^2 = 1 \\
\alpha x + \beta y + \gamma z = 0
\end{cases}.
\end{align*}
计算$\int_C y(y - x)\mathrm{d}s$.
\end{example}
\begin{solution}

记$K\triangleq \sqrt{\alpha^2+\beta^2}$，$M\triangleq \sqrt{\alpha^2+\beta^2+\gamma^2}$，考虑如下正交变换
\begin{align*}
\begin{pmatrix}
x'\\
y'\\
z'
\end{pmatrix}
=
\begin{pmatrix}
\frac{\alpha\gamma}{KM} & \frac{\beta\gamma}{KM} & -\frac{K}{M}\\
-\frac{\beta}{K} & \frac{\alpha}{K} & 0\\
\frac{\alpha}{M} & \frac{\beta}{M} & \frac{\gamma}{M}
\end{pmatrix}
\begin{pmatrix}
x\\
y\\
z
\end{pmatrix}
= T
\begin{pmatrix}
x\\
y\\
z
\end{pmatrix}.
\end{align*}
于是曲线$C$方程等价于
\begin{align*}
\begin{cases}
\left( \frac{\alpha\gamma}{KM}x'-\frac{\beta}{K}y' \right)^2+\left( \frac{\beta\gamma}{KM}x'+\frac{\alpha}{K}y' \right)^2+\left( -\frac{K}{M}x' \right)^2=1\\
z' = 0
\end{cases}
\Longleftrightarrow
\begin{cases}
x'^2+y'^2=1\\
z' = 0
\end{cases}.
\end{align*}
因为正交变换是等距变换，进而保持弧长微元，即
\begin{align*}
\mathrm{d}s=\sqrt{(\mathrm{d}x)^2+(\mathrm{d}y)^2+(\mathrm{d}z)^2}
= \left| \begin{pmatrix} \mathrm{d}x\\ \mathrm{d}y\\ \mathrm{d}z \end{pmatrix} \right|
= \left| T\begin{pmatrix} \mathrm{d}x'\\ \mathrm{d}y'\\ \mathrm{d}z' \end{pmatrix} \right|
= \left| \begin{pmatrix} \mathrm{d}x'\\ \mathrm{d}y'\\ \mathrm{d}z' \end{pmatrix} \right|
= \sqrt{(\mathrm{d}x')^2+(\mathrm{d}y')^2+(\mathrm{d}z')^2}
= \mathrm{d}s'.
\end{align*}
又由曲线$C:\begin{cases}x'^2+y'^2=1\\z'=0\end{cases}$的对称性知
\begin{align*}
\int_{x'^2+y'^2=1,z'=0}x'y'\mathrm{d}s'
= \int_{x'^2+y'^2=1,z'=0}(-x')y'\mathrm{d}s'
\Longrightarrow
\int_{x'^2+y'^2=1,z'=0}x'y'\mathrm{d}s'=0.
\end{align*}
所以
\begin{align*}
\int_C y(y-x)\mathrm{d}s
=& \int_{x'^2+y'^2=1,z'=0}\left( \frac{\beta\gamma}{KM}x'+\frac{\alpha}{K}y' \right)\left( \frac{(\beta-\alpha)\gamma}{KM}x'+\frac{\alpha+\beta}{K}y' \right)\mathrm{d}s'\\
=& \int_{x'^2+y'^2=1,z'=0}\left[ \frac{\beta(\beta-\alpha)\gamma^2}{K^2M^2}x'^2+\frac{\alpha(\alpha+\beta)}{K^2}y'^2 \right]\mathrm{d}s'\\
=& \int_0^{2\pi}\left[ \frac{\beta(\beta-\alpha)\gamma^2}{K^2M^2}\cos^2\theta+\frac{\alpha(\alpha+\beta)}{K^2}\sin^2\theta \right]\mathrm{d}\theta\\
=& \frac{\beta(\beta-\alpha)\gamma^2}{K^2M^2}\pi + \frac{\alpha(\alpha+\beta)}{K^2}\pi
= \pi \frac{\alpha^2+\alpha\beta+\gamma^2}{\alpha^2+\beta^2+\gamma^2}.
\end{align*}
\end{solution}


























































































\end{document}